\section{The carry-free source and the digit operator}
\label{sec:digit}

\subsection{Conventions and phase removal}

For \(d\ge1\), define
\[
  C_d=(\one_{\{a+c<d\}})_{0\le a,c<d},
  \qquad
  \kappa_{d,q}=\|C_d\|_{S_q},
\]
and let \(C_0\) denote the zero block, with \(\kappa_{0,q}=0\).  The
positive-vertex alphabet is always \(\N=\{1,2,\ldots\}\).  When a
length-\(L\) tensor \(C_b^{\otimes L}\) is used as a finite control, its
all-zero word is therefore not a vertex of the infinite source.

The imaginary part of \(s\) has no effect on singular values, but the
relation is a two-sided unitary equivalence rather than a unitary
conjugation.  Indeed, for \(s=\sigma+\ii t\), let
\[
  U_te_n=n^{-\ii t/2}e_n.
\]
Entrywise,
\begin{equation}
  B_{b,s}=U_t B_{b,\sigma}U_t.
  \label{eq:two-sided-unitary}
\end{equation}
Both factors are unitary, so every singular-value assertion can be proved
for the real symmetric matrix \(B_{b,\sigma}\).  We do not infer from
\cref{eq:two-sided-unitary} that \(B_{b,s}\) is positive or self-adjoint
when \(t\ne0\).

\subsection{Exact singular values}

\begin{proposition}[Digit spectrum]
\label{prop:digit-spectrum}
The singular values of \(C_b\), in decreasing order, are
\begin{equation}
  s_j(C_b)=
  [
    2\sin\frac{(2j-1)\pi}{4b+2}
  ]^{-1},
  \qquad 1\le j\le b.
  \label{eq:digit-spectrum}
\end{equation}
Consequently,
\begin{equation}
  \kappa_{b,q}
  =
  [
    \sum_{j=1}^{b}
    (
      2\sin\frac{(2j-1)\pi}{4b+2}
    )^{-q}
  ]^{1/q}.
  \label{eq:kappa-explicit}
\end{equation}
\end{proposition}

\begin{proof}
The \((a,a')\) entry of \(C_bC_b^*\) counts the digits \(c\) satisfying
\(c<b-\max(a,a')\), and hence equals \(b-\max(a,a')\).  Reverse the row
and column order and index the result by \(1,\ldots,b\).  The resulting
matrix is
\[
  K_b=(\min(i,j))_{1\le i,j\le b}.
\]
Its inverse is the tridiagonal matrix
\[
  L_b=K_b^{-1}
  =
  \begin{pmatrix}
    2&-1\\
    -1&2&-1\\
    &\ddots&\ddots&\ddots\\
    &&-1&2&-1\\
    &&&-1&1
  \end{pmatrix}.
\]
For completeness, multiplying \(K_b\) by this matrix subtracts consecutive
columns twice in the interior and once at the terminal boundary, producing
the identity.

The eigenvalue equation for \(L_b\) has the interior solution
\(v_i=\sin(i\theta)\), with
\(\lambda=2-2\cos\theta=4\sin^2(\theta/2)\) and the left boundary
\(v_0=0\).  The final row is equivalent to the ghost boundary
\(v_{b+1}=v_b\), so
\[
  \sin((b+1)\theta)-\sin(b\theta)=0.
\]
The \(b\) roots in \((0,\pi)\) are
\(\theta_j=(2j-1)\pi/(2b+1)\).  Thus the eigenvalues of \(K_b\) are
\([4\sin^2(\theta_j/2)]^{-1}\).  Taking square roots proves
\cref{eq:digit-spectrum}, and summing their \(q\)th powers proves
\cref{eq:kappa-explicit}.
\end{proof}

\begin{corollary}[Two exact comparisons]
\label{cor:digit-comparisons}
For every \(b\ge2\),
\begin{equation}
  \kappa_{b,2}^2=\frac{b(b+1)}2<b^2,
  \qquad
  \tau_b:=\kappa_{b,1}>b.
  \label{eq:digit-comparisons}
\end{equation}
In particular \(\log_b\kappa_{b,2}<1\) and
\(\alpha_b:=\log_b\tau_b>1\).
\end{corollary}

\begin{proof}
The Hilbert--Schmidt norm squared counts the ones in \(C_b\), giving
\(b+(b-1)+\cdots+1=b(b+1)/2\).  Reversing the columns of \(C_b\) makes it
triangular with unit diagonal, so
\(\lvert\det C_b\rvert=1\).  The product of its singular values is
therefore one.  Their arithmetic mean is at least their geometric mean,
and the inequality is strict because their squared sum
\(b(b+1)/2\) exceeds \(b\).  Hence \(\tau_b>b\).
\end{proof}

\begin{table}[t]
  \centering
  \small
  % Generated by scripts/generate_paper_assets.py; do not edit by hand.
\begin{tabular}{@{}rrrrrr@{}}
\toprule
$b$ & $\tau_b$ & $\alpha_b$ & $\kappa_{b,2}$ & $\log_b\kappa_{b,2}$ & $\sigma_c(q{=}2)$ \\
\midrule
2 & 2.236068 & 1.160964 & 1.732051 & 0.792481 & 1.000000 \\
3 & 3.603875 & 1.166936 & 2.449490 & 0.815465 & 1.000000 \\
4 & 5.064178 & 1.170164 & 3.162278 & 0.830482 & 1.000000 \\
5 & 6.595934 & 1.172119 & 3.872983 & 0.841303 & 1.000000 \\
\bottomrule
\end{tabular}

  \caption{Display evaluations of the exact formulas in
  \cref{eq:kappa-explicit,eq:digit-comparisons}.  The last column is
  \(\sigma_c(q=2)=\max\{1,\log_b\kappa_{b,2}\}\); the decimals are not
  empirical estimates and do not define the theorem.}
  \label{tab:digit-thresholds}
\end{table}

At \(b=2\), \cref{eq:digit-spectrum} gives the golden-ratio pair
\((1+\sqrt5)/2\) and \((\sqrt5-1)/2\).  Their sum is \(\sqrt5\), so
\begin{equation}
  \tau_2=\sqrt5,
  \qquad
  \alpha_2=\log_2\sqrt5.
  \label{eq:binary-alpha}
\end{equation}
These identities will distinguish the binary trace-class wall from its
Hilbert--Schmidt wall.
